% =======================================================================================
% MODELO DE TCC - INSTITUTO FEDERAL DE SERGIPE (IFS)
% =======================================================================================
% Autores do Modelo (discentes de Ciência da Computação - IFS, Campus Itabaiana):
%   Mikael dos Santos Gois
%   Talyson Rodrigues Silva Nascimento
%   Kauan Gutierrez Santos Estácio
% Modelo Base: abnTeX2, adaptado para as normativas internas de TCC do IFS (Art. 31 e 34)
% Créditos dos modelos base (abnTeX2, DECOM/UFOP e Diogo D. Moreira) e das imagens: ver NOTICE
% Licença: CC BY-SA 4.0 (ver LICENSE)
% Este arquivo faz parte de um modelo opensource para estudantes do IFS.
% Não remova os créditos dos autores originais e do modelo base.
% Este modelo é fornecido "no estado em que se encontra", sem garantias de qualquer tipo.
% =======================================================================================

% [O que aparece no PDF Bookmark] {O que aparece impresso no documento}
\chapter[Trabalhos Correlatos]{TRABALHOS CORRELATOS} \label{sec:trabalhos-correlatos}

A análise de trabalhos correlatos permite mapear o estado da arte e identificar as abordagens metodológicas adotadas por outros pesquisadores para resolver problemas semelhantes. Nesta seção, são discutidas publicações recentes relacionadas à execução de algoritmos de inferência e coleta de dados inerciais.

Ao longo da literatura, nota-se que grande parte das implementações prioriza a execução de modelos computacionais em plataformas de alto desempenho. O Quadro \ref{qua:comparativo_trabalhos} sintetiza as principais diferenças de arquitetura e hardware entre os trabalhos analisados e a solução proposta nesta pesquisa.

% ---------------------------------------------------------------------------------------
% EXEMPLO DE QUADRO (DADOS TEXTUAIS) - NBR 14724:2024, 4.2.1.9 e 5.8
% ---------------------------------------------------------------------------------------
% DICA: dados TEXTUAIS (comparações, características, definições) vão em QUADRO. A norma
% trata o quadro como ILUSTRAÇÃO: identificação com palavra designativa + número +
% travessão + título, fonte logo abaixo e entrada na LISTA DE QUADROS (elemento próprio).
% No quadro, a moldura fechada é permitida.
% DICA: cite o quadro no texto ANTES dele aparecer (com \ref{}) e substitua o conteúdo
% de exemplo pelo seu.
\begin{quadro}[htb]
    \centering
    \caption{Comparativo entre os trabalhos correlatos e a solução proposta}
    \vspace{0.2cm}
    
    % DICA PARA O ESTUDANTE:
    % 'l' = alinhado à esquerda (ideal para colunas curtas, como o nome do Autor)
    % 'c' = centralizado (ideal para colunas com respostas curtas, como Sim/Não)
    % 'X' = coluna dinâmica que quebra o texto automaticamente e preenche o espaço da página
    \begin{tabularx}{\textwidth}{|l | X | X | X|}
        \hline
        \textbf{Autor (Ano)} & \textbf{[Inserir Critério 1]} & \textbf{[Inserir Critério 2]} & \textbf{[Inserir Critério 3]} \\
        \hline
        Sobrenome \emph{et al.} (2023) & Método longo que agora vai quebrar a linha sozinho & Baixo & Não \\
        \hline
        Sobrenome e Silva (2024)& Método Y & Médio & Sim \\
        \hline
        Sobrenome (2025)        & Método Z & Alto  & Não \\
        \hline
        \textbf{Proposta Atual} & \textbf{Método Inovador} & \textbf{Baixo} & \textbf{Sim} \\
        \hline
    \end{tabularx}
    
    \fonte{Elaborada pelo autor (2026).}
    \label{qua:comparativo_trabalhos}
\end{quadro}

Conforme evidenciado no Quadro \ref{qua:comparativo_trabalhos}, a principal distinção da proposta atual reside no.... Enquanto os trabalhos de Sobrenome \emph{et al.} (2023) e Sobrenome e Silva (2024) apresentam limitações quanto a [Citar o problema deles], a presente abordagem inova ao utilizar [Sua solução]. Essa escolha garante [Citar vantagem], preenchendo a lacuna identificada na literatura analisada.

% ---------------------------------------------------------------------------------------
% EXEMPLO DE TABELA (DADOS NUMÉRICOS) - NBR 14724:2024, 5.9 + Normas do IBGE
% ---------------------------------------------------------------------------------------
% DICA: dados NUMÉRICOS vão em TABELA, padronizada pelas Normas de apresentação tabular
% do IBGE (a NBR 14724:2024, 5.9, remete a elas):
%   - a moldura NÃO deve ter traços verticais que delimitem a tabela à esquerda e à
%     direita (IBGE, 4.3.3) - por isso o tabular abaixo não usa "|" nas laterais;
%   - no MÍNIMO três traços horizontais: topo, cabeçalho e rodapé (IBGE, 4.3.1);
%   - traços verticais apenas para DESTACAR parte do cabeçalho ou dos dados (IBGE, 4.3.2);
%   - a unidade de medida deve ser indicada no cabeçalho da coluna (IBGE, 4.6);
%   - recomenda-se alinhar os dados numéricos à direita (boa prática de leitura; a
%     norma do IBGE não fixa o alinhamento).
% DICA: se a tabela não couber em uma página, repita o cabeçalho e indique "continua"
% (primeira página), "continuação" (intermediárias) e "conclusão" (última), deixando o
% rodapé apenas na página de conclusão (IBGE).
% DICA: a tabela entra na LISTA DE TABELAS e é citada no texto com \ref{}.
A Tabela \ref{tab:selecao_rsl} exemplifica um dado numérico: a quantidade de trabalhos recuperados e selecionados em cada base de dados consultada.

\begin{table}[htb]
    \centering
    \caption{Trabalhos recuperados e selecionados por base de dados}
    \label{tab:selecao_rsl}
    \begin{tabular}{lrrr}
        \hline
        \textbf{Base de dados} & \textbf{Recuperados} & \textbf{Selecionados} & \textbf{Aproveitamento (\%)} \\
        \hline
        ACM Digital Library & 412 & 38 & 9,2 \\
        IEEE Xplore         & 356 & 41 & 11,5 \\
        Scopus              & 289 & 27 & 9,3 \\
        \textbf{Total}      & \textbf{1\,057} & \textbf{106} & \textbf{10,0} \\
        \hline
    \end{tabular}
    \fonte{Elaborada pelo autor (2026).}
\end{table}

Os números da Tabela \ref{tab:selecao_rsl} mostram que [explique, com suas palavras, o que os dados indicam: quantos trabalhos foram recuperados, quantos passaram pelos critérios de seleção e o que o percentual de aproveitamento revela sobre a busca].

% DICA PARA O ESTUDANTE:
% Sempre comente os resultados da tabela (ou do quadro) no parágrafo imediatamente
% abaixo, explicando com suas palavras o que os dados mostram.
%
% DICA: a expressão latina "et al." é grafada EM ITÁLICO no texto, nas citações
% e nas referências - assim como a própria ABNT a grafia nos exemplos da
% NBR 10520:2023 (6.1.2) e da NBR 6023:2025 (8.1.1.2). Use \emph{et al.}
% em textos escritos à mão; nas citações automáticas (\cite) e nas
% referências o próprio modelo já gera o itálico (entrada @abnt-options
% "opcoes-abnt" em referencias.bib + \citeoption{opcoes-abnt} em main.tex).